%! Multiplying \& Dividing Radicals; Rationalizing — Unit 6, Lesson 6.3 — Warm-Up
\documentclass[10pt]{article}
\usepackage{algebra2-article}
\usepackage{algebra2-boxes}
\newcommand{\TallMath}[1]{$\displaystyle #1\rule[-1.4em]{0pt}{3.2em}$}

\begin{document}

\pageheader{Unit 6, Lesson 6.3}{Warm-Up}
\namedateperiod

\vspace{0.10in}

\begin{notesbox}{Getting Ready: Skills We'll Use Today}
\small Three skills you already have. Today they turn into one new one. Work quickly.

\vspace{4pt}
\textbf{1. Distributing and multiplying binomials.} Expand and simplify.
\begin{center}
\renewcommand{\arraystretch}{1.6}
\begin{tabular}{@{}l l l@{}}
(a) \; $3x(2x+5)=$ \blank{2.6cm} & (b) \; $(x+4)(x-4)=$ \blank{2.4cm} &
(c) \; $(x+4)^{2}=$ \blank{2.8cm} \\
\end{tabular}
\end{center}
In (b) two of the four terms cancelled and the answer had no middle term. What was special about
those two binomials? \writeline

\vspace{4pt}
\textbf{2. Yesterday's radicals.} Simplify each.
\begin{center}
\renewcommand{\arraystretch}{1.6}
\begin{tabular}{@{}l l l l@{}}
(a) \; $\sqrt{12}=$ \blank{1.4cm} & (b) \; $\sqrt{50}=$ \blank{1.4cm} &
(c) \; $\sqrt{7}\cdot\sqrt{7}=$ \blank{1.4cm} & (d) \; $\left(\sqrt{11}\right)^{2}=$ \blank{1.4cm} \\
\end{tabular}
\end{center}
Look at (c) and (d). What does multiplying a square root \emph{by itself} do to the radical sign?
\writeline

\vspace{4pt}
\textbf{3. Multiplying by a form of $1$.}
\begin{itemize}[leftmargin=1.5em, itemsep=3pt]
  \item Fill in the missing numerator: \; $\dfrac{3}{4}=\dfrac{\blank{1.0cm}}{20}$. \;
        What did you multiply the top \emph{and} the bottom by? \blank{1.2cm}
  \item Multiplying a fraction by $\dfrac{5}{5}$ does not change its value, because
        $\dfrac{5}{5}=\blank{1.0cm}$. \; Is the same true of $\dfrac{\sqrt{2}}{\sqrt{2}}$?
        \blank{1.4cm}
\end{itemize}

\end{notesbox}

\end{document}
